\documentclass[11pt]{article}
\usepackage[margin=1in]{geometry}
\usepackage{amsmath,amssymb,amsthm,mathtools}
\usepackage{bm}
\usepackage{enumitem}
\usepackage{hyperref}
\usepackage{xcolor}

\hypersetup{colorlinks=true,linkcolor=blue,citecolor=blue,urlcolor=blue}

\title{Derivation of the Finite-Dimensional Model and Discussion of Actuator Placement}
\author{Generated from Project Handout}
\date{\today}

\newcommand{\R}{\mathbb{R}}
\newcommand{\Om}{\Omega}
\newcommand{\dd}{\,d}
\newcommand{\p}{\partial}
\newcommand{\ip}[2]{\left\langle #1,#2\right\rangle}
\newtheorem*{remark}{Remark}

\begin{document}
% \maketitle

\section{Derivation of the Finite-Dimensional Model}

We start with the controlled wave equation for the membrane displacement $u(x,y,t)$:
\begin{equation}
\label{eq:wave}
 u_{tt}(x,y,t) = c^2 \Delta u(x,y,t) + b(t)\,\psi(x,y), 
 \qquad (x,y)\in\Omega,\ t>0,
\end{equation}
where $\Omega = (0,1)\times(0,1)$ is the unit square, with boundary conditions
\begin{equation}
\label{eq:bc}
 u(x,y,t)=0,\qquad (x,y)\in \partial\Omega,
\end{equation}
and initial conditions
\begin{equation}
\label{eq:ic}
 u(x,y,0)=u_0(x,y),\qquad u_t(x,y,0)=v_0(x,y).
\end{equation}

The Dirichlet Laplacian eigenfunctions on the unit square are
\begin{equation}
\label{eq:eigs}
\phi_{mn}(x,y)=2\sin(m\pi x)\sin(n\pi y),\qquad m,n\ge 1,
\end{equation}
with eigenvalues
\begin{equation}
\label{eq:eigvals}
\lambda_{mn}=\pi^2(m^2+n^2).
\end{equation}
These form an orthonormal basis of $L^2(\Omega)$.

We expand the solution as
\begin{equation}
\label{eq:expansion}
 u(x,y,t)=\sum_{m=1}^{\infty}\sum_{n=1}^{\infty} q_{mn}(t)\,\phi_{mn}(x,y).
\end{equation}

Substituting into the PDE \eqref{eq:wave} and projecting onto mode $(m,n)$ yields
\begin{equation}
\label{eq:modalode}
 \ddot q_{mn}(t) + c^2\lambda_{mn} q_{mn}(t) = \beta_{mn}\, b(t),
\end{equation}
where
\begin{equation}
\label{eq:beta-general}
\beta_{mn}=\int_0^1\!\int_0^1 \psi(x,y)\phi_{mn}(x,y)\,\dd x\,\dd y.
\end{equation}
For a point actuator at $(x_0,y_0)$, this simplifies to
\begin{equation}
\label{eq:beta-point}
\beta_{mn}=\phi_{mn}(x_0,y_0)=2\sin(m\pi x_0)\sin(n\pi y_0).
\end{equation}

To obtain a finite-dimensional approximation, we truncate to modes with $1 \leq m,n \leq M$, giving $N = M^2$ modes. We flatten the indices and define the state vector
\[
q(t)=\begin{bmatrix} q_1(t) \\ \vdots \\ q_N(t) \end{bmatrix},
\qquad
p(t)=\dot q(t).
\]
Let
\[
\omega_k^2=c^2\lambda_k,
\qquad
\beta=\begin{bmatrix}\beta_1 \\ \vdots \\ \beta_N\end{bmatrix},
\qquad
\Omega_N = \operatorname{diag}(\omega_1^2,\dots,\omega_N^2).
\]
The system becomes
\begin{equation}
\label{eq:firstorder}
\dot x = A x + B b,
\qquad
x=\begin{bmatrix} q \\ p \end{bmatrix}\in\R^{2N},
\end{equation}
with
\begin{equation}
\label{eq:AB}
A=
\begin{bmatrix}
0 & I \\
-\Omega_N & 0
\end{bmatrix},
\qquad
B=
\begin{bmatrix}
0 \\
\beta
\end{bmatrix}.
\end{equation}

The energy of the system is
\begin{equation}
\label{eq:energy}
E(t)=\frac12\sum_{k=1}^N \left(p_k(t)^2 + \omega_k^2 q_k(t)^2\right).
\end{equation}

For LQR control, we minimize the cost
\begin{equation}
\label{eq:cost}
J(b)=\int_0^{\infty} \left( x(t)^T Q x(t) + R\,b(t)^2 \right)\dd t,
\end{equation}
with a physically motivated choice
\begin{equation}
\label{eq:Qchoice}
Q=
\begin{bmatrix}
\alpha\,\Omega_N & 0 \\
0 & \beta_v I
\end{bmatrix},
\qquad \alpha>0,\ \beta_v>0.
\end{equation}

The optimal feedback is $b(t) = -K x(t)$, where $K = R^{-1} B^T P$ and $P$ solves the algebraic Riccati equation
\begin{equation}
\label{eq:ARE}
A^T P + P A - P B R^{-1} B^T P + Q = 0.
\end{equation}

\section{Discussion of Actuator Placement and Uncontrollable Modes}

The actuator location significantly affects which modes can be directly controlled. The coupling coefficient for mode $(m,n)$ is
\[
\beta_{mn} = 2 \sin(m\pi x_0) \sin(n\pi y_0)
\]
for the point-actuator model.

If $\beta_{mn} = 0$ for a particular mode, that mode is not directly actuated and may remain uncontrollable. This is particularly problematic for symmetric placements.

For example, at the center $(x_0, y_0) = (0.5, 0.5)$,
\[
\beta_{mn} = 2 \sin\left(\frac{m\pi}{2}\right) \sin\left(\frac{n\pi}{2}\right),
\]
which vanishes whenever $m$ or $n$ is even. Thus, many modes are missed, potentially leaving the system underactuated.

To ensure better controllability, an off-center location such as $(x_0, y_0) = (0.37, 0.61)$ is recommended, as it breaks symmetry and typically provides nonzero coupling to the first many modes.

In practice, actuator placement should be chosen to maximize the number of controllable modes within the truncation, ensuring effective stabilization of the membrane vibrations.

\end{document}